\documentclass[openright, oneside, a4paper, article, brazil]{abntex2}
\usepackage{formularios}
% --------- %
% Metadados %
% --------- %
\hypersetup{%
	colorlinks=false,
	allbordercolors=white,
	pdftitle = {Análise — Isenção de anuidade por questões de saúde},
	pdfauthor = {Conselho Regional de Psicologia de São Paulo},
	pdfsubject = {Isenção de anuidade},
	pdfkeywords = {%
		Isenção de anuidade,
		saúde,
	}%
}
\newlength{\nomeda}
	\settowidth{\nomeda}{Nome da/o psicóloga/o: }
\newlength{\nome}
	\setlength{\nome}{\textwidth-\nomeda-0.5em}
\title{Análise — Isenção de anuidade por questões de saúde}
\author{Conselho Regional de Psicologia de São Paulo}
% --------- %
% DOCUMENTO %
% --------- %
\begin{document}
\pagestyle{crp}
\centering
\Form
\section*{ANÁLISE --- ISENÇÃO DE ANUIDADE POR QUESTÕES DE SAÚDE}%
% -------------------------- %
% Identificação/qualificação %
% -------------------------- %
\vfill
\setasuspacing{\OnehalfSpacing}
% Define espaçamento duplo para que as caixas de texto não fiquem sobrepostas
\justifying
\preenctext{14}{nome}{\nome}{Nome da/o psicóloga/o: \hfill}
\vspace{\baselineskip}

\preenctext{14}{crp}{8em}{Nº de inscrição no CRP~SP: } \hfill Data de inscrição: \dataabrev{14}{1}
\vspace{\baselineskip}

Situação da inscrição:\hspace{1em} \caixasel{14}{ativa} ativa;\hspace{1em} \caixasel{14}{cancelada} cancelada.
\vfill

\justifying
\subsection{Solicita isenção das anuidades}

\begin{enumerate}
	\item \preenctext{14}{anuidade1}{47em}{\hfill}
	\item \preenctext{14}{anuidade2}{47em}{\hfill}
	\item \preenctext{14}{anuidade3}{47em}{\hfill}
	\item \preenctext{14}{anuidade4}{47em}{\hfill}
	\item \preenctext{14}{anuidade5}{47em}{\hfill}
\end{enumerate}
\vfill

\subsection{Documentos recebidos}

\begin{enumerate}
	\item \preenctext{14}{documento1}{47em}{\hfill}
	\item \preenctext{14}{documento2}{47em}{\hfill}
	\item \preenctext{14}{documento3}{47em}{\hfill}
	\item \preenctext{14}{documento4}{47em}{\hfill}
	\item \preenctext{14}{documento5}{47em}{\hfill}
\end{enumerate}
\vfill

\setasuspacing{\OnehalfSpace}\justifying
\begin{framed}
	\noindent\small Obs.: de acordo com o \href{https://transparencia.cfp.org.br/wp-content/uploads/sites/7/2020/05/Manual-de-Procedimentos-Administrativos-e-Financeiro-CFP-RES-20.2018.pdf}{\textcolor{crp1}{\textit{Manual de procedimentos administrativos, financeiros e contábeis do Sistema Conselhos de Psicologia}}}, ``A Diretoria do Conselho Regional deverá decidir em 30 (trinta) dias, cabendo recurso ao Plenário, no prazo de 20 (vinte) dias, se indeferido o pedido [\textit{de interrupção temporária do pagamento das anuidades}]. Não havendo deliberação pela Diretoria, decorridos 30 (trinta) dias do recebimento do pedido, a interrupção temporária será considerada aprovada'' (Norma 02, seção 5.4.2).
\end{framed}
\vfill
\clearpage
\setasuspacing{\DoubleSpace}\justifying
\subsection{Parecer técnico}
\preenctext{14}{parecer1}{\textwidth}{\null}
\preenctext{14}{parecer2}{\textwidth}{\null}
\preenctext{14}{parecer3}{\textwidth}{\null}
\preenctext{14}{parecer4}{\textwidth}{\null}
\preenctext{14}{parecer5}{\textwidth}{\null}
%\vspace{0.5\baselineskip}

\begin{framed}\noindent
	\begin{minipage}{0.5\linewidth}{%
			\subsection{Processo ético (PE)}
			\begin{minipage}{\boxlarg}
				\caixasel{14}{pe-sim}
			\end{minipage}%
			\begin{minipage}{\linewidth-\boxlarg}
				Sim
			\end{minipage}
			\begin{minipage}{\boxlarg}
				\caixasel{14}{pe-nao}
			\end{minipage}%
			\begin{minipage}{\linewidth-\boxlarg}
				Não
			\end{minipage}
		}%
	\end{minipage}%
	\begin{minipage}{0.5\linewidth}{%
		\subsection{Providências}
		\begin{minipage}{\boxlarg}
			\caixasel{14}{deferir}
		\end{minipage}%
		\begin{minipage}{\linewidth-\boxlarg}
			Deferir
		\end{minipage}
		\begin{minipage}{\boxlarg}
			\caixasel{14}{indeferir}
		\end{minipage}%	
		\begin{minipage}{\linewidth-\boxlarg}
			Indeferir
		\end{minipage}
	}%
	\end{minipage}
\end{framed}

\subsection{Deliberação}

\preenctext{14}{deliberacao1}{\textwidth}{\null}
\preenctext{14}{deliberacao2}{\textwidth}{\null}
\preenctext{14}{deliberacao3}{\textwidth}{\null}
\preenctext{14}{deliberacao4}{\textwidth}{\null}
\preenctext{14}{deliberacao5}{\textwidth}{\null}

\begin{flushright}
	\preenctext{14}{responsavel-analise}{24em}{\null}\\
	\footnotesize Responsável pela análise
\end{flushright}
\vspace*{4\baselineskip}
\begin{center}
	\rule{24em}{0.4pt}\\
	\footnotesize Conselheira/o secretária/o
\end{center}

\end{document}